%! AP Statistics — Unit 8, Lesson 8.6 — Homework
\documentclass[10pt]{article}
\usepackage{apstats-article}
\usepackage{apstats-boxes}
\usepackage{grading}
\gradingassignment{apstats-u08-l06-hw}
\gradingcourse{advanced_statistics}
\newcommand{\TallMath}[1]{$\displaystyle #1\rule[-1.4em]{0pt}{3.2em}$}

\begin{document}

\pageheader{Unit 8, Lesson 8.6}{Homework}

\vspace{0.08in}

\textit{Show all work. Use all four SPDC steps for Problem 1.}

\begin{enumerate}[itemsep=10pt]

%% Problem 1 %%
\item \textbf{Sleep and Grade Level.}
A school counselor randomly selects 180 teenagers and classifies them by
\textbf{grade} (9th / 10th / 11th) and \textbf{sleep amount} ($<7$ hours / $7{+}$ hours per night).

\renewcommand{\arraystretch}{1.4}
\begin{center}
\begin{tabular}{@{} l c c c @{}}
\hline
           & \textbf{$<$7 hrs} & \textbf{7$+$ hrs} & \textbf{Total} \\
\hline
9th grade  & 30 & 30 & 60  \\
10th grade & 38 & 22 & 60  \\
11th grade & 42 & 18 & 60  \\
\hline
\textbf{Total} & 110 & 70 & 180 \\
\hline
\end{tabular}
\end{center}

A chi-square table excerpt (right-tail areas):

\renewcommand{\arraystretch}{1.3}
\begin{center}\small
\begin{tabular}{@{} c c c c c @{}}
\hline
\textbf{df} & $p = 0.10$ & $p = 0.05$ & $p = 0.025$ & $p = 0.01$ \\
\hline
1  & 2.706  & 3.841  & 5.024  & 6.635  \\
2  & 4.605  & 5.991  & 7.378  & 9.210  \\
3  & 6.251  & 7.815  & 9.348  & 11.345 \\
4  & 7.779  & 9.488  & 11.143 & 13.277 \\
\hline
\end{tabular}
\end{center}

\begin{enumerate}[label=\alph*., itemsep=6pt]
  \boxguard[18]
  \item \textbf{State.} Identify the test type. Write $H_0$ and $H_a$.

        Test type: \gradedhide{\writeline\\\gradedhere{\gradedblank[id=hw1a1, type=rubric, points=1, display={\bfseries Independence}]{\linewidth+6pt}{\gradedlinesheight{9}}}\writelines{10}}

        \boxguard[14]
        $H_0$: \gradedhide{\writeline\\\gradedhere{\gradedblank[id=hw1a2, type=rubric, points=1, display={\bfseries No association.}]{\linewidth+6pt}{\gradedlinesheight{9}}}\writelines{10}}

        \boxguard[14]
        $H_a$: \gradedhide{\writeline\\\gradedhere{\gradedblank[id=hw1a3, type=rubric, points=1, display={\bfseries Association.}]{\linewidth+6pt}{\gradedlinesheight{9}}}\writelines{10}}

  \boxguard[16]
  \item \textbf{Plan.} Method: \gradedhide{\writeline\\\gradedhere{\gradedblank[id=hw1b1, type=rubric, points=1, display={\bfseries $\chi^2$ test (same as a.)}]{\linewidth+6pt}{\gradedlinesheight{9}}}\writelines{10}}

        \boxguard[14]
        Verify the Random condition: \gradedhide{\writeline\\\gradedhere{\gradedblank[id=hw1b2, type=rubric, points=1, display={\bfseries Stated.}]{\linewidth+6pt}{\gradedlinesheight{9}}}\writelines{10}}

        Compute expected counts (each row total = 60):

        $E_{\text{9th, }<7} = \dfrac{60 \times 110}{180} = $ \makebox[2.5cm][r]{\smash{\gradedblank[id=hw1b3, type=short, points=1, answer=36.67, check=numeric, tolerance=0.01]{0.95cm}{0.2in}}}
        \quad $E_{\text{9th, }7+} = \dfrac{60 \times 70}{180} = $ \hphantom{\blank{2.5cm}}\hskip 0pt plus 1fill\smash{\gradedblank[id=hw1b4, type=short, points=1, answer=23.33, check=numeric, tolerance=0.01]{1.6cm}{0.2in}}

        Are all expected counts $\ge 5$? \hphantom{\blank{2cm}}\hskip 0pt plus 1fill\smash{\gradedblank[id=hw1b5, type=short, points=1, answer=yes, check=text]{1.2cm}{0.2in}}

  \newpage
  \item \textbf{Do.} $\text{df} = ($ \blank{1.5cm} $-1)($ \blank{1.5cm} $-1) = $ \hphantom{\blank{2cm}}\hskip 0pt plus 1fill\smash{\gradedblank[id=hw1c1, type=short, points=1, answer=2, check=numeric, tolerance=0]{1.2cm}{0.2in}}

        Compute $\chi^2$ contributions and sum them below.

        \renewcommand{\arraystretch}{1.6}
        \begin{center}
        \begin{tabular}{@{} l c c @{}}
        \hline
                   & \textbf{$<$7 hrs} & \textbf{7$+$ hrs} \\
        \hline
        9th grade  & $\dfrac{(30 - \blank{2cm})^2}{\blank{2cm}} = $ \blank{2cm}
                   & $\dfrac{(30 - \blank{2cm})^2}{\blank{2cm}} = $ \blank{2cm} \\
        10th grade & \blank{4.5cm} & \blank{4.5cm} \\
        11th grade & \blank{4.5cm} & \blank{4.5cm} \\
        \hline
        \end{tabular}
        \end{center}

        $\chi^2 = $ \makebox[4cm][r]{\smash{\gradedblank[id=hw1c2, type=short, points=2, answer=5.25, check=numeric, tolerance=0.1]{2.2cm}{0.2in}}}

        \boxguard[14]
        Using the table above, the p-value is approximately \gradedhide{\writeline\\\gradedhere{\gradedblank[id=hw1c3, type=rubric, points=1, display={\bfseries 0.05 to 0.10}]{\linewidth+6pt}{\gradedlinesheight{9}}}\writelines{10}}

  \item \textbf{Conclude.} At $\alpha = 0.05$:
        \gradedhide{\writeline\\\gradedhere{\gradedblank[id=hw1d, type=rubric, points=2, display={$p \approx 0.07 > 0.05$: fail to reject $H_0$; no convincing evidence of an association between grade and sleep.}]{\linewidth+6pt}{\gradedlinesheight{1}}}\writelines{9}}
\end{enumerate}

\vspace{0.1in}

\newpage
%% Problem 2 %%
\item \textbf{Degrees of Freedom.}
A student computes $\chi^2 = 9.21$ for a $2 \times 4$ contingency table.

\begin{enumerate}[label=\alph*.]
  \item What is the degrees of freedom? \makebox[3cm][r]{\smash{\gradedblank[id=hw2a, type=short, points=1, answer=3, check=numeric, tolerance=0]{1.45cm}{0.2in}}}
  \item Using the table in Problem 1, what is the critical value at $\alpha = 0.05$?
        \makebox[3cm][r]{\smash{\gradedblank[id=hw2b, type=short, points=1, answer=7.815, check=numeric, tolerance=0]{1.45cm}{0.2in}}}
  \boxguard[16]
  \item Based on $\chi^2 = 9.21$, can the student reject $H_0$ at $\alpha = 0.05$?
        Explain.
        \gradedhide{\writeline\\\gradedhere{\gradedblank[id=hw2c, type=rubric, points=2, display={\bfseries Yes: $9.21 > 7.815$.}]{\linewidth+6pt}{\gradedlinesheight{9}}}\writelines{9}}
\end{enumerate}

\vspace{0.1in}

%% Problem 3 %%
\item \textbf{Conceptual Question.}
In what sense does a \emph{larger} $\chi^2$ value provide more evidence against $H_0$?
Use the ideas of expected counts, p-value, and the chi-square distribution in your answer.
\gradedhide{\writeline\\\gradedhere{\gradedblank[id=hw3, type=rubric, points=3, display={Larger $\chi^2$ = observed counts farther from expected under $H_0$; farther into the right tail, so a smaller p-value and stronger evidence against $H_0$.}]{\linewidth+6pt}{\gradedlinesheight{3}}}\writelines{9}}

\end{enumerate}

\end{document}
